\documentclass{article}
\usepackage[T1]{fontenc}
\usepackage{lmodern}
\usepackage{graphicx}
\usepackage{amsmath,amssymb}
\usepackage[a4paper,margin=2cm]{geometry}
\usepackage[absolute,overlay]{textpos}
\setlength{\TPHorizModule}{1mm}
\setlength{\TPVertModule}{1mm}
\usepackage[indonesian]{babel}
\usepackage{hyperref}
\setlength{\emergencystretch}{2em}
% unit: Halaman 14

\begin{document}

% === Halaman 14 (llm) ===

\begin{itemize}
    \item Contoh di dalam Vigenere Cipher, cipherteks setiap huruf plainteks dihitung dengan:
    \[ C = (P + K) \pmod{26} \]
    \item Penyerang memilih plaintext P = AAAAAAA... (= 000000...).
    \item Maka cipherteks yang keluar, C, pasti sama dengan K $(C = K)$.
    \item Dengan mengetahui K, cipherteks lain bisa langsung dicoba didekripsi dengan K tersebut.
\end{itemize}

\end{document}
